% === Begin Label Data ===
% ID: 415
% 难度星级: 
% 标签: 
% 备注: 
% 组卷引用次数: 0
% === End  Label Data ===

\begin{problem}{2021}{G}{新高考I卷}{16}{排列组合}
某校学生在研究民间剪纸艺术时，发现剪纸时经常会沿纸的某条对称轴把纸对折。规格为 $ 20\,\mathrm{dm}\times 12\,\mathrm{dm} $ 的长方形纸，对折 $ 1 $ 次共可以得到 $ 10\,\mathrm{dm}\times 12\,\mathrm{dm} $、$ 20\,\mathrm{dm}\times 6\,\mathrm{dm} $ 两种规格的图形，它们的面积之和 $ S_1=240\,\mathrm{dm}^2 $；对折 $ 2 $ 次共可以得到 $ 5\,\mathrm{dm}\times 12\,\mathrm{dm} $、$ 10\,\mathrm{dm}\times 6\,\mathrm{dm} $、$ 20\,\mathrm{dm}\times 3\,\mathrm{dm} $ 三种规格的图形，它们的面积之和 $ S_2=180\,\mathrm{dm}^2 $；以此类推。则对折 $ 4 $ 次共可以得到不同规格图形的种数为 \underline{\hspace{4em}}；如果对折 $ n $ 次，那么 $ \displaystyle\sum_{k=1}^{n}S_k= $ \underline{\hspace{4em}}$ \mathrm{dm}^2 $.
\end{problem}

\begin{answer}

\end{answer}

\begin{solutions}

\end{solutions}